% ==== $X$ ====
\yoschapter{X}{Compact Operators}
\label{chap:X}

Let $X$ and $Y$ be complex $B$-spaces, and let $S$ be the unit sphere in $X$. An operator $T \in L(X, Y)$ is said to be compact or completely continuous if the image $T \cdot S$ is relatively compact in $Y$. For a compact operator $T \in L(X, X)$, the eigenvalue problem can be treated fairly completely, in the sense that the classical theory of Fredholm concerning linear integral equations may be extended to the linear functional equation $Tx - \lambda x = y$ with a complex parameter $\lambda$. This result is known as the Riesz-Schauder theory. F. \textsc{Riesz} \cite{rickart273} and J. \textsc{Schauder} \cite{schauder289}. 


\yossection{1}{Compact Sets in \texorpdfstring{$B$}{B}-spaces}
\label{sec:X.1}

A compact set in a linear topological space must be bounded. The converse is, however, not true in general; we know (Chapter~\ref{chap:III}, \ref{sec:III.2}) that the closed unit sphere of a normed linear space $X$ is strongly compact iff $X$ is of finite dimension. Let $S$ be a compact metric space and $C(S)$ the $B$-space of real- or complex-valued continuous functions $x(s)$ on $S$, normed by $\|x\| = \sup\limits_{s \in S} |x(s)|$. We know (Chapter~\ref{chap:III}, \ref{sec:III.3}) that a subset $\{x_{\alpha}(s)\}$ of $C(S)$ is strongly relatively compact in $C(S)$ iff $\{x_{\alpha}(s)\}$ is equi-bounded and equi-continuous in $\alpha$. For the case of the space $L^p(S, \mathfrak{B}, m)$, $1 \leqq p < \infty$, we have 

\begin{theoremn}[\textsc{Fréchet}-\textsc{Kolmogorov}] Let $S$ be the real line, $\mathfrak{B}$ the $\sigma$-ring of Baire subsets $B$ of $S$ and $m(B) = \int_{B} \dd x$ the ordinary Lebesgue measure of $B$. Then a subset $K$ of $L^{p}(S, \mathfrak{B}, m)$, $1 \leqq p < \infty$, is strongly pre-compact iff it satisfies the conditions:
\begin{align}
\sup\limits _{x \in K} \| x \|& = \sup\limits _{x \in K} \left(\int_{S} | x (s) | ^{p} \dd s\right) ^{1/p} < \infty , \label{eq:X.1.1} \\
\lim\limits _{t \to 0}& \int_{S} | x (t + s) - x (s) | ^{p} \dd s = 0 \text{ uniformly in } x \in K, \label{eq:X.1.2} \\
\lim\limits _{\alpha \uparrow \infty}& \int_{| s | > \alpha} | x (s) | ^{p} \dd s = 0 \text{ uniformly in } x \in K. \label{eq:X.1.3}
\end{align}

\end{theoremn}
\begin{proof}
 Let $K$ be strongly relatively compact. Then $K$ is bounded and so \eqref{eq:X.1.1} is true. Let $\varepsilon > 0$ be given. Then there exists a finite number of functions $\in L^p: f_1, f_2, \ldots, f_n$ such that, for each $f \in K$, there is a $j$ with $\|f - f_j\| \leqq \varepsilon$. Otherwise, we would have an infinite sequence $\{f_j\} \subseteq K$ with $\|f_j - f_i\| > \varepsilon$ for $j \neq i$, contrary to the relative compactness of $K$. We then find, by the definition of the Lebesgue integral, finitely-valued functions $g_1, g_2, \ldots, g_n$ such that $\|f_j - g_j\| \leqq \varepsilon$ ( $j = 1, 2, \ldots, n$). Since each finitely-valued function $g_j(x)$ vanishes outside some sufficiently large interval, we have, for large $\alpha$, 
\[
\begin{gathered}
\left(\int_{\alpha} ^{\infty} + \int_{- \infty} ^{- \alpha} | f (s) | ^{p} \dd s\right) ^{1/p} \leqq \left(\int_{\alpha} ^{\infty} + \int_{- \infty} ^{- \alpha} | f (s) - g _{j} (s) | ^{p} \dd s\right) ^{1/p}
\\
+ \left(\int_{\alpha} ^{\infty} + \int_{- \infty} ^{- \alpha} | g _{j} (s) | ^{p} \dd s\right) ^{1/p} \leqq \| f - g _{j} \| + \left(\int_{\alpha} ^{\infty} + \int_{- \infty} ^{- \alpha} | g _{j} (s) | ^{p} \dd s\right) ^{1/p}. \end{gathered}
\]

This proves \eqref{eq:X.1.3} by $\| f - g_j\| \leqq \| f - f_j\| + \| f_j - g_j\| \leqq 2\varepsilon$. 

The proof of \eqref{eq:X.1.2} is based on the fact that, for the defining function $C_I(s)$ of a finite interval $I$, $\lim\limits_{t \to 0} \int_{-\infty}^{\infty} |C_I(s + t) - C_I(s)|^p \dd s = 0$ (see Chapter~\ref{chap:0}, \ref{sec:0.3}). Thus \eqref{eq:X.1.2} holds for finitely-valued functions $g_j(s)$ ( $j = 1, 2, \ldots, n$). Hence we have, for any $f \in K$, 
\[
\begin{aligned}
&\overline {{\lim\limits _{t \to 0}}} \left(\int_{- \infty} ^{\infty} | f (s + t) - f (s) | ^{p} \dd s\right) ^{1/p} \leqq {} \overline {{\lim\limits _{t \to 0}}} \left(\int_{- \infty} ^{\infty} | f (s + t) - f _{j} (s + t) | ^{p} \dd s\right) ^{1/p}
\\
&\quad + \overline {{\lim\limits _{t \to 0}}} \left(\int_{- \infty} ^{\infty} | f _{j} (s + t) - g _{j} (s + t) | ^{p} \dd s\right) ^{1/p} + \overline {{\lim\limits _{t \to 0}}} \left(\int_{- \infty} ^{\infty} | g _{j} (s + t) - g _{j} (s) | ^{p} \dd s\right) ^{1/p}
\\
&\quad + \left(\int_{- \infty} ^{\infty} | g _{j} (s) - f _{j} (s) | ^{p} \dd s\right) ^{1/p} + \left(\int_{- \infty} ^{\infty} | f _{j} (s) - f (s) | ^{p} \dd s\right) ^{1/p} \leqq {} \varepsilon + \varepsilon + 0 + \varepsilon + \varepsilon ,
\end{aligned}
\]
by taking $f_{j}$ in such a way that $\| f - f_{j}\| \leqq \varepsilon$. This proves \eqref{eq:X.1.2}. 
\end{proof}

\begin{theoremn}[We next prove the converse part of the Theorem] We define the translation operator $T_{t}$ by $(T_{t}f)(s) = f(t + s)$. Condition \eqref{eq:X.1.2} says that  
$s\text{- }\lim\limits_{t \to 0} T_t f = f$ uniformly in $f \in K$. We next define the mean value $(M_a f)(s) = (2a)^{-1} \int_{-a}^{a} (T_t f)(s) \dd t$. Then, by Hölder's inequality and the Fubini-Tonelli theorem, 
\[
\begin{aligned}
\left\| M _{a} f - f \right\|& \leqq \left(\int_{- \infty} ^{\infty} \left\{\int_{-a} ^{a} (2 a) ^{-1} | f (s + t) - f (s) | \dd t \right\} ^{p} \dd s\right) ^{1/p}
\\
\quad& \leqq (2 a) ^{-1} \left(\int_{- \infty} ^{\infty} \int_{-a} ^{a} | f (s + t) - f (s) | ^{p} \dd t \cdot (2 a) ^ {p / p ^{\prime}} \dd s\right) ^{1/p}
\\
\quad& \leqq \left((2 a) ^{-1} \int_{-a} ^{a} \dd t \int_{- \infty} ^{\infty} | f (s + t) - f (s) | ^{p} \dd s\right) ^{1/p} \text{ if } 1 \leqq p < \infty .
\end{aligned}
\]
\end{theoremn}

Thus we have 
\[
\left\| M _{a} f - f \right\| \leqq \sup\limits _{| t | \leqq a} \left\| T _{i} f - f \right\|,
\]
so that $s\text{- }\lim\limits_{a \downarrow 0} M_a f = f$ uniformly in $f \in K$. Therefore, we have to prove the relative compactness of the set $\{M_a f; f \in K\}$ for a sufficiently small fixed $a > 0$. 

We will show that, for a fixed $a > 0$, the set of functions $\{(M_{a}f)(s); f \in K\}$ is equi-bounded and equi-continuous. In fact, we have, as above, 
\[
\begin{aligned} \left| (M _{a} f) (s _{1}) - (M _{a} f) (s _{2}) \right| & \leqq (2 a) ^{-1} \int_{-a} ^{a} | f (s _{1} + t) - f (s _{2} + t) | \dd t \\ & \leqq \left((2 a) ^{-1} \int_{-a} ^{a} | f (s _{1} + t) - f (s _{2} + t) | ^{p} \dd t\right) ^{1/p}. \end{aligned}
\]

Thus, by \eqref{eq:X.1.2}, we have proved the equi-continuity of the set of functions $\{(M_{a}f)(s); f \in K\}$ for a fixed $a > 0$. The equi-boundedness of the set may be proved similarly. Thus, by the Ascoli-Arzelà theorem, there exists, for any positive $a > 0$, a finite number of functions $M_{a}f_{1}, M_{a}f_{2}, \ldots, M_{a}f_{n}$ with $f_{j} \in K$ ( $j = 1, 2, \ldots, n$) such that, for any $f \in K$, there exists some $j$ for which $\sup |(M_{a}f)(s) - (M_{a}f_{j})(s)| \leqq \varepsilon$. Therefore 
\begin{equation}\label{eq:X.1.4}
\begin{aligned}
&\left\| M _{a} f - M _{a} f _{j} \right\| ^{p}  \leqq \int_{- \alpha} ^{\alpha} | (M _{a} f) (s) - (M _{a} f _{j}) (s) | ^{p} \dd s + \int_{| s | > \alpha} | (M _{a} f) (s) - (M _{a} f _{j}) (s) | ^{p} \dd s.
\end{aligned}
\end{equation}

The second term on the right is, by Minkowski's inequality, smaller than $\left(\|M_a f - f\| + \left(\int_{|s| > \alpha} |f(s) - f_j(s)|^p \dd s\right)^{1/p} + \left(\int_{|s| > \alpha} |f_j(s) - (M_a f_j)(s)|^p \dd s\right)^{1/p}\right)^p$. 

The term $\|M_{a}f - f\|$ is small for sufficiently small $a > 0$, and, by virtue of \eqref{eq:X.1.3}, $\int_{|s| > \alpha}|f(s) - f_j(s)|^p \dd s$ and $\int_{|s| > \alpha}|f_j(s) - (M_{a}f_j)(s)|^p \dd s$ are both small for sufficiently large $\alpha > 0$, if $a > 0$ is bounded. Also the first term on the right of \eqref{eq:X.1.4} is $\leqq 2\alpha\varepsilon^{p}$ for an appropriate choice of $j$. These estimates are valid uniformly with respect to $f \in K$. Thus we have proved the relative compactness in $L^{p}$ of the set $\{M_{a}f; f \in K\}$ for sufficiently small $a > 0$.\yosstarfootnote{See also Supplementary Notes, p.~\pageref{notes:X}.} 


\yossection{2}{Compact Operators and Nuclear Operators}
\label{sec:X.2}

\begin{definition}\label{def:X.2.Definition.1} Let $X$ and $Y$ be $B$-spaces, and let $S$ be the unit sphere of $X$. An operator $T \in L(X, Y)$ is said to be compact or completely continuous if the image $T \cdot S$ is relatively compact in $Y$. 
\end{definition}

\begin{example}\label{ex:X.2.Example.1} Let $K(x, y)$ be a real- or complex-valued continuous function defined for $-\infty < a \leqq x, y \leqq b < \infty$. Then the integral operator $K$ defined by 
\begin{equation}\label{eq:X.2.1}
(K f) (x) = \int_{a} ^{b} K (x, y) f (y) \dd y
\end{equation}
is compact as an operator $\in L(C[a,b],C[a,b])$. 
\end{example}

\begin{proof}
 Clearly $K$ maps $C[a, b]$ into $C[a, b]$. Set $\sup |K(x, y)| = M$. 

Then $\| K \cdot f \| \leqq (b - a) M \| f \|$ so that $K \cdot S$ is equi-bounded. By Schwarz' inequality, we have 
\[
\left| (K f) \left(x _{1}\right) - (K f) \left(x _{2}\right) \right| ^{2} \leqq \int_{a} ^{b} \left| K \left(x _{1}, y\right) - K \left(x _{2}, y\right) \right| ^{2} \dd y \cdot \int_{a} ^{b} | f (y) | ^{2} \dd y,
\]
and hence $K \cdot S$ is equi-continuous, that is, 
\[
\lim\limits _{\delta \downarrow 0} \sup\limits _ {| x _{1} - x _{2} | \leqq \delta} | (K f) (x _{1}) - (K f) (x _{2}) | = 0 \text{ uniformly in } f \in S.
\]

Therefore, by the Ascoli-Arzelà theorem (Chapter~\ref{chap:III}, \ref{sec:III.3}), the set $K \cdot S$ is relatively compact in $C[a, b]$. 
\end{proof}

\begin{example}\label{ex:X.2.Example.2} Let $K(x, y)$ be a real- or complex-valued $\mathfrak{B}$-measurable function on a measure space $(S, \mathfrak{B}, m)$ such that 
\begin{equation}\label{eq:X.2.2}
\iint_{S} \left| K (x, y) \right| ^{2} m (\dd x) m (\dd y) < \infty .
\end{equation}
\end{example}

Then the integral operator $K$ defined by the kernel $K(x, y)$: 
\begin{equation}\label{eq:X.2.3}
(K f) (x) = \int_{S} K (x, y) f (y) m (\dd y), f \in L ^{2} (S) = L ^{2} (S, \mathfrak {B}, m),
\end{equation}
is compact as an operator $\in L(L^2(S), L^2(S))$. The kernel $K(x, y)$ satisfying \eqref{eq:X.2.2} is said to be of the Hilbert-Schmidt type. 

\begin{proof}
 Take any sequence $\{f_n\}$ from the unit sphere of $L^2(S)$. We have to show that the sequence $\{K \cdot f_n\}$ is relatively compact in $L^2(S)$. Since a Hilbert space $L^2(S)$ is locally sequentially weakly compact, we may assume that $\{f_n\}$ converges weakly to an element $f \in L^2(S)$; otherwise, we choose a suitable subsequence. By \eqref{eq:X.2.2} and the Fubini-Tonelli theorem, $\int_{S} |K(x, y)|^2 m(\dd y) < \infty$ for $m$-a.e. $x$. Hence, for such an $x$, $\lim\limits_{n \to \infty} (Kf_n)(x) = \lim\limits_{n \to \infty} \int_{S} K(x, y) f_n(y) m(\dd y) = \lim\limits_{n \to \infty} (f_n(.), \overline{K(x,.)}) = (f(.), \overline{K(x,.)}) = \int_{S} K(x, y) f(y) m(\dd y)$. On the other hand, we have, by Schwarz' inequality, 
\begin{equation}\label{eq:X.2.4}
\begin{aligned}
&\left| (K f _{n}) (x) \right| ^{2} \leqq \int_{S} | K (x, y) | ^{2} m (\dd y) \cdot \int_{S} \left| f _{n} (y) \right| ^{2} m (\dd y) \leqq \int_{S} | K (x, y) | ^{2} m (\dd y) \text{ for } m\text{-a.e.} x.
\end{aligned}
\end{equation}

Hence, by the Lebesgue-Fatou theorem, $\lim\limits_{n\to \infty}\int_S|(Kf_n)(x)|^2 m(\dd x) = \int_S|(Kf)(x)|^2 m(\dd x)$. This result, if combined with $w - \lim\limits_{n\to \infty}K\cdot f_n = K\cdot f$, implies $\slim\limits_{n\to \infty}K\cdot f_n = K\cdot f$ by Theorem~\ref{thm:V.1.Theorem.8} in Chapter~\ref{chap:V}, \ref{sec:V.1}. But, as proved above in \eqref{eq:X.2.4}, we have 
\[
\int_{S} | (K h) (x) | ^{2} m (\dd x) \leqq \iint_{S} | K (x, y) | ^{2} m (\dd y) m (\dd x) \cdot \int_{S} | h (y) | ^{2} m (\dd y),
\]
and so 
\begin{equation}\label{eq:X.2.5}
\left\| K \right\| \leqq \left(\iint_{S} \left| K (x, y) \right| ^{2} m (\dd x) m (\dd y)\right) ^{1/2}.
\end{equation}

Hence, from \wlim $f_n = f$, we obtain \wlim $K \cdot f_n = K \cdot f$, because, for any $g \in L^2(S)$, $\lim\limits_{n \to \infty} (K \cdot f_n, g) = \lim\limits_{n \to \infty} (f_n, K^*g) = (f, K^*g) = (K \cdot f, g)$. 
\end{proof}

\begin{theoremu}
 \begin{enumerate}[(i)]
\item A linear combination of compact operators is compact.
\item The product of a compact operator with a bounded linear operator is compact; thus the set of compact operators $\in L(X, X)$ constitutes a closed two-sided ideal of the algebra $L(X, X)$ of operators.
\item Let a sequence $\{T_n\}$ of compact operators $\in L(X, Y)$ converge to an operator $T$ in the sense of the uniform operator topology, i.e., $\lim\limits_{n \to \infty} \|T - T_n\| = 0$. Then $T$ is also compact.
\end{enumerate}
\end{theoremu}

\begin{proof}
 (i) and (ii) are clear from the definition of compact operators. The closedness, in the sense of the uniform operator topology, of the ideal of compact operators in the algebra $L(X, X)$ is implied by (iii). 

We shall prove (iii). Let $\{x_h\}$ be a sequence from the closed unit sphere $S$ of $X$. By the compact property of each $T_n$, we can choose, by the diagonal method, a subsequence $\{x_{h'}\}$ such that $s\lim\limits_{h\to \infty}T_n x_{h'}$ exists for every fixed $n$. We have 
\[
\begin{aligned} & {\left\| T x _{h ^{\prime}} - T x _{k ^{\prime}} \right\| \leqq \left\| T x _{h ^{\prime}} - T _{n} x _{h ^{\prime}} \right\| + \left\| T _{n} x _{h ^{\prime}} - T _{n} x _{k ^{\prime}} \right\| + \left\| T _{n} x _{k ^{\prime}} - T x _{k ^{\prime}} \right\|} \\ & {\quad \leqq \| T - T _{n} \| + \left\| T _{n} x _{h ^{\prime}} - T _{n} x _{k ^{\prime}} \right\| + \| T _{n} - T \|,} \end{aligned}
\]
and so $\varlimsup\limits_{h,k\to\infty}\left\|T\cdot x_{h'}-T\cdot x_{k'}\right\|\leqq 2\|T-T_n\|$. Hence $\{Tx_{k'}\}$ is a Cauchy sequence in the $B$-space $Y$. 
\end{proof}

Nuclear Operator. As an application of the Theorem, we shall consider the nuclear operator introduced by A. \textsc{Grothendieck} \cite{grothendieck88}. 

\begin{definition}\label{def:X.2.Definition.2} Let $X, Y$ be $B$-spaces and $T \in L(X, Y)$. If there exist a sequence $\{f_n'\} \subseteq X'$, a sequence $\{y_n\} \subseteq Y$ and a sequence $\{c_n\}$ of numbers such that 
\begin{align}
\sup\limits _{n} \left\| f _{n} ^{\prime} \right\|& < \infty , \sup\limits _{n} \left\| y _{n} \right\| < \infty , \sum\limits_{n} \left| c _{n} \right| < \infty \text{ and } \notag \\
T \cdot x& = \slim _{m \rightarrow \infty} \sum\limits_{n=1} ^{m} c _{n} \left\langle x, f _{n} ^{\prime} \right\rangle y _{n} \text{ in } Y \text{ for every } x \in X, \label{eq:X.2.6}
\end{align}
then $T$ is called a nuclear operator on $X$ into $Y$. 
\end{definition}

\begin{remarku} The existence of the \slim in \eqref{eq:X.2.6} is clear, since 
\[
\left\| \sum\limits_{j=n} ^{m} c _{j} \langle x, f _{j} ^{\prime} \rangle y _{j} \right\| \leqq \sum\limits_{j=n} ^{m} | c _{j} | \cdot \| x \| \cdot \| f _{j} ^{\prime} \| \cdot \| y _{j} \| \leqq \text{ constant } \cdot \sum\limits_{j=n} ^{m} | c _{j} | \cdot \| x \|.
\]
\end{remarku}

The nuclear condition says that the \slim is equal to $T \cdot x$ for every $x \in X$. 

\begin{propositionu} $A$ nuclear operator $T$ is compact. 
\end{propositionu}

\begin{proof}
 Define the operator $T_{n}$ by 
\begin{equation}\label{eq:X.2.7}
T _{n} x = \sum\limits_{j=1} ^{n} c _{j} \langle x, f _{j} ^{\prime} \rangle y _{j}. \end{equation}

Since the range $R(T_{n})$ is of finite dimension, $T_{n}$ is compact as may be proved by the Bolzano-Weierstrass theorem. Moreover, by \eqref{eq:X.2.6} and 
\[
\left\| T x - T _{n} x \right\| = \left\| \sum\limits_{j=n+1} ^{\infty} c _{j} \langle x, f _{j} ^{\prime} \rangle y _{j} \right\| \leqq \text{ constant } \sum\limits_{j=n+1} ^{\infty} \left| c _{j} \right| \cdot \left\| x \right\|,
\]
we have $\lim\limits_{n\to \infty}\left\|T - T_n\right| = 0$ and so $T$ must be compact. 
\end{proof}

An Example of the Nuclear Operator. Let $G$ be a bounded open domain of $\mathbb{R}^{n}$, and consider the Hilbert space $H_{0}^{k}(G)$. Suppose $(k-j)>n$. Then the mapping $T$ 
\begin{equation}\label{eq:X.2.8}
H _{0} ^{k} (G) \ni \varphi \rightarrow \varphi \in H _{0} ^{j} (G)
\end{equation}
is a nuclear operator $\in L(H_0^k (G),H_0^j (G))$ 

\begin{proof}
 We may assume that the bounded domain $G$ is contained in the interior of the parallelogram $P$: 
\[
0 \leqq x _{j} \leqq 2 \pi (j = 1, 2, \dots , n). \]

We recall that $H_0^k (G)$ is the completion of $\hat{H}_0^k (G) = C_0^k (G)$ with respect to the norm $\|\varphi \|_k = \left(\sum\limits_{|s|\leqq k}\int_G |D^s\varphi (x)|^2 \dd x\right)^{1/2}$ (see Chapter~\ref{chap:I}, \ref{sec:I.10}). We 
extend the functions $\in \tilde{H}_0^k (G)$ to be periodic with period $2\pi$ in each variable $x_{s}$ by defining the function values as $0$ in $P - G$. The functions 
% 按正文段落排: 左对齐到版心、首行缩进, 编号排在末行右端。
% 用 \notag 关掉居中编号, 改在段内用 \hfill 把 (\theequation) 推到右边,
% 这样公式体就是整版心宽的段落, 不再是居中排的块。
\begin{equation}\notag
\refstepcounter{equation}\label{eq:X.2.9}%
\parbox[b]{\linewidth}{\setlength{\parindent}{1.2em}\setlength{\parfillskip}{0pt}$f _{\beta} (x) = (2 \pi) ^{-n/2} \exp (i \beta \cdot x)$, where $\beta = (\beta_{1}, \beta_{2}, \dots , \beta_{n})$ is an $n$-tuple of integers and \mbox{$\beta \cdot x = \sum\limits_{s=1}^{n} \beta_s x_s,$}\hfill(\theequation)}
\end{equation}


form a complete orthonormal system of $L^2(P) = H_0^0(P)$. Thus, denoting by $D^s$ the distributional derivative, we have, for $|s| \leqq k$, the Fourier expansion in $L^2(P)$ of functions $D^s\varphi(x)$ where $\varphi \in \hat{H}_0^k(G)$: 
\begin{equation}\label{eq:X.2.10}
D ^{s} \varphi (x) = \sum\limits_{\beta} \left(D ^{s} \varphi , f _{\beta}\right) _{0} f _{\beta}, \text{ where } (\psi , f _{\beta}) _{0} = \int_{P} \psi (x) \overline {{f _{\beta} (x)}} \dd x.
\end{equation}

We have, by 
\[
(D ^{s} \varphi , f _{\beta}) _{0} = (- 1) ^{|s|} (\varphi , D ^{s} f _{\beta}) _{0} = \prod\limits_{m=1} ^{n} (i \beta_{m}) ^{sm} (\varphi , f _{\beta}) _{0}
\]
and the Parseval relation 
\[
\sum\limits_{\beta} \left| (D ^{s} \varphi , f _{\beta}) _{0} \right| ^{2} = \int_{P} \left| D ^{s} \varphi (x) \right| ^{2} \dd x \leqq \left\| \varphi \right\| _{k} ^{2} (| s | \leqq k),
\]
the inequality 
\[
\left| \left(\varphi , (1 + | \beta | ^{2}) ^{k/2} f _{\beta}\right) _{0} \right| ^{2} \leqq \text{ constant } \sum\limits_{| s | \leqq k} \left| \left(D ^{s} \varphi , f _{\beta}\right) _{0} \right| ^{2} \leqq \text{ constant } \| \varphi \| _{k} ^{2}.
\]

Therefore the functional $f_{\beta}^{\prime} \in H_0^k (G)'$ defined by 
\[
\langle \varphi , f _{\beta} ^{\prime} \rangle = (\varphi , (1 + | \beta | ^{2}) ^{k/2} f _{\beta}) _{0}
\]
satisfies $\sup\limits_{\beta}\|f_{\beta}'\| < \infty$. Moreover, 
\[
y _{\beta} = (1 + | \beta | ^{2}) ^{-j/2} f _{\beta}
\]
satisfies $\sup\limits_{\beta}\|y_{\beta}\|_j < \infty$ by $D^{s}f_{\beta} = \prod\limits_{t=1}^{n}(i\beta_{t})^{s_{t}}f_{\beta}$. We also have 
\[
\sum\limits_{\beta} | c _{\beta} | < \infty , \text{ where } c _{\beta} = (1 + | \beta | ^{2}) ^{(j-k)/2},
\]
because, for positive integers $\beta_{s}$, 
\[
\begin{aligned}
\sum\limits_{\beta} \frac {1}{(\beta_{1} + \beta_{2} + \cdots + \beta_{n}) ^{k-j}} &= \sum\limits_{\beta} \left(\frac {1}{(\beta_{1} + \beta_{2} + \cdots + \beta_{n}) ^{n}}\right) ^{(k-j)/n}
\\
\leqq \sum\limits_{\beta} \left(\frac {1}{\beta_{1} \beta_{2} \cdots \beta_{n}}\right) ^{(k-j)/n} &= \sum\limits_ {\beta_{1}} \left(\frac {1}{\beta_{1}}\right) ^{(k-j)/n} \cdot \sum\limits_ {\beta_{2}} \left(\frac {1}{\beta_{2}}\right) ^{(k-j)/n}
\\
\dots \sum\limits_ {\beta_{n}} \left(\frac {1}{\beta_{n}}\right) ^{(k-j)/n}& < \infty \text{ by } \frac {(k - j)}{n} > 1.
\end{aligned}
\]

Therefore, we have proved the (Fourier) expansion 
\[
\varphi = \sum\limits_{\beta} c _{\beta} \left< \varphi , f _{\beta} ^{\prime} \right> y _{\beta}.
\]
\end{proof}

\begin{remarku} If there exists a complete orthonormal system $\{\varphi_{j}\}$ of eigenfunctions of a given bounded linear operator $K \in L(L^{2}(S), L^{2}(S))$ such that $K\varphi_{j} = \lambda_{j}\varphi_{j} (j = 1, 2, \ldots)$, then, from the Fourier expansion 
\[
f = \sum\limits_{j=1} ^{\infty} (f, \varphi_{j}) \varphi_{j}, f \in L ^{2} (S),
\]
\end{remarku}

we obtain 
\[
K f = \sum\limits_{j=1} ^{\infty} \lambda_{j} (f, \varphi_{j}) \varphi_{j}.
\]

We have $\lambda_{j} = (K\varphi_{j},\varphi_{j})$, and so, if the eigenvalues $\lambda_{j}$ are all $>0$ and $\sum\limits_{j=1}^{\infty}\lambda_{j}<\infty$, then the operator $K$ is nuclear. The condition $\sum\limits_{j=1}^{\infty}|(K\varphi_{j},\varphi_{j})|<\infty$ is surely satisfied, if the operator $K$ is defined by a kernel 
\[
K(x,y)=\int_S\overline{K_2(z,x)}K_1(z,y)m(\dd z),
\]
where the kernels $K_1(x,y)$ and $K_2(x,y)$ are of the Hilbert-Schmidt type.
For,
\[
\sum\limits_{j=1} ^{\infty} \left| \left(K _{2} ^{*} K _{1} \varphi_{j}, \varphi_{j}\right) \right| = \sum\limits_{j=1} ^{\infty} \left| \left(K _{1} \varphi_{j}, K _{2} \varphi_{j}\right) \right| \leqq \left(\sum\limits_{j=1} ^{\infty} \left\| K _{1} \varphi_{j} \right\| ^{2} \cdot \sum\limits_{j=1} ^{\infty} \left\| K _{2} \varphi_{j} \right\| ^{2}\right) ^{1/2}
\]
and, by Parseval's relation, we have 
\[
\begin{aligned}
&\sum\limits_{j=1} ^{\infty} \| K _{1} \varphi_{j} \| ^{2}  = \sum\limits_{j=1} ^{\infty} \int_{S} \left| \int_{S} K _{1} (z, y) \varphi_{j} (y) m (\dd y) \right| ^{2} m (\dd z) \\
&= \int_{S} \sum\limits_{j=1} ^{\infty} \left| \int_{S} K _{1} (z, y) \varphi_{j} (y) m (\dd y) \right| ^{2} m (\dd z) = \int_{S} \left\{\int_{S} | K _{1} (z, y) | ^{2} m (\dd y) \right\} m (\dd z) < \infty ,
\end{aligned}
\]
and similar for $\sum\limits_{j=1}^{\infty} \|K_2\varphi_j\|^2$. A bounded linear operator $K$ in a separable Hilbert space $X$ is said to be of the trace class if $\sum\limits_{j=1}^{\infty} |(K\varphi_j, \psi_j)| < \infty$ for arbitrary complete orthonormal systems $\{\varphi_j\}$ and $\{\psi_j\}$ of $X$. For a general account concerning the trace class and the nuclear operator, see R. \textsc{Schatten} \cite{schatten288}, and I. M. \textsc{Gelfand}-N. Y. \textsc{Vilenkin} \cite{gelfand77}. 


\yossection{3}{The Rellich-Gårding Theorem}
\label{sec:X.3}

\begin{theoremn}[\textsc{Garding} {\cite{garding74}}] Let $G$ be a bounded open domain of $\mathbb{R}^{n}$. If an operator $T \in L(H_0^k(G), H_0^k(G))$ satisfies, for $j < k$,
\begin{equation}\label{eq:X.3.1}
\|T\varphi\|_{k} \leqq C \|\varphi\|_{j} \quad \text{for all } \varphi \in H_{0}^{k}(G) \text{, where } C \text{ is a constant,}
\end{equation}
then $T$ is compact as an operator $\in L(H_{0}^{k}(G), H_{0}^{k}(G))$. 

\end{theoremn}
\begin{proof}
 By the definition of the space $H_0^k(G)$ (see Chapter~\ref{chap:I}, \ref{sec:I.10}), it would be sufficient to show the following: Let a sequence $\{\varphi_\nu\} \subseteq \hat{H}_0^k(G) = C_0^k(G)$ be such that $\|\varphi_\nu\|_k \leqq 1 (\nu = 1,2,\ldots)$. Then the sequence $\{T\varphi_\nu\}$ contains a subsequence strongly convergent in $H_0^k(G)$. The Fourier transform $\hat{\varphi}_\nu(\xi) = (2\pi)^{-n/2} \int_G \varphi_\nu(x) \exp(-ix\xi) \dd x$ satisfies, by Schwarz' inequality, 
\[
\left| \hat {\varphi} _{\nu} (\xi) \right| ^{2} \leqq (2 \pi) ^{-n} \int_{G} \dd x \int_{G} \left| \varphi_{\nu} (x) \right| ^{2} \dd x \leqq (2 \pi) ^{-n} \int_{G} \dd x,
\]
and hence $\{\hat{\varphi}_{\nu}(\xi)\}$ is equi-bounded in $\xi\in \mathbb{R}^{n}$ and in $\pmb{\nu}$. We may assume, by the boundedness of $\|\varphi_{\nu}\|_{0}$, that a subsequence $\{\varphi_{\nu}\}$ is weakly convergent in $L^2(G) = H_0^0(G)$. Since, for each $\xi$, the function $\exp(-ix\xi)$ belongs to $L^2(G)$, we know that the sequence of bounded functions $\hat{\varphi}_{\nu'}(\xi) = (\varphi_{\nu'}, (2\pi)^{-n/2} \exp(-ix\xi))_0$ converges at every $\xi$. Thus, by \eqref{eq:X.3.1} and the Parseval relation for the Fourier transform (Chapter~\ref{chap:VI}, \ref{sec:VI.2}), 
\[
\begin{aligned}
&\left\| T \varphi_ {\nu^{\prime}} - T \varphi_ {\mu^{\prime}} \right\| _{k} ^{2} = \left\| T (\varphi_ {\nu^{\prime}} - \varphi_ {\mu^{\prime}}) \right\| _{k} ^{2} \leqq C ^{2} \left\| \varphi_ {\nu^{\prime}} - \varphi_ {\mu^{\prime}} \right\| _{j} ^{2}
\\
&= C ^{2} \sum\limits_{| s | \leqq j} \left\| D ^{s} (\varphi_ {\nu^{\prime}} - \varphi_ {\mu^{\prime}}) \right\| _{0} ^{2} = C ^{2} \sum\limits_{| s | \leqq j} \left\| (D ^{s} (\varphi_ {\nu^{\prime}} - \varphi_ {\mu^{\prime}})) \right\| _{0} ^{2}
\\
&= C ^{2} \sum\limits_{| s | \leqq j} \left\| \prod\limits_{t=1} ^{n} (i \xi_{t}) ^{st} (\hat {\varphi} _ {\nu^{\prime}} - \hat {\varphi} _ {\mu^{\prime}}) (\xi) \right\| _{0} ^{2} \leqq C ^{2} \sum\limits_{| s | \leqq j} \int_{| \xi | \leqq r} \left| \prod\limits_{t=1} ^{n} \xi_{t} ^{st} (\hat {\varphi} _ {\nu^{\prime}} (\xi) - \hat {\varphi} _ {\mu^{\prime}} (\xi)) \right| ^{2} \dd \xi
\\
&\quad + C ^{2} C _{1} \int_{| \xi | > r} | \xi | ^{2j} | \hat {\varphi} _ {\nu^{\prime}} (\xi) - \hat {\varphi} _ {\mu^{\prime}} (\xi) | ^{2} \dd \xi ,
\end{aligned}
\]
where $C_1$ is a positive constant. 

The first term on the right converges, for fixed $r$, to $0$ as $v'$ and $\mu' \to \infty$. This we see by the Lebesgue-Fatou Lemma. The second term on the right is, for $r > 1$. 
\[
\begin{aligned}
&\leqq C ^{2} C _{1} r ^{2j-2k} \int_{| \xi | > r} | \xi | ^{2k} | \hat {\varphi} _ {\nu^{\prime}} (\xi) - \hat {\varphi} _ {\mu^{\prime}} (\xi) | ^{2} \dd \xi \leqq C ^{2} C _{1} r ^{2j-2k} \int_ {\mathbb{R}^{n}} | \xi | ^{2k} | \hat {\varphi} _ {\nu^{\prime}} (\xi) - \hat {\varphi} _ {\mu^{\prime}} (\xi) | ^{2} \dd \xi
\\
&\leqq C ^{2} C _{1} C _{2} r ^{2j-2k} \sum\limits_{| s | \leqq k} \| (D ^{s} \widehat {\varphi_ {\nu^{\prime}} - D ^{s}} \varphi_ {\mu^{\prime}}) \| _{0} ^{2} = C ^{2} C _{1} C _{2} r ^{2j-2k} \sum\limits_{| s | \leqq k} \| D ^{s} (\varphi_ {\nu^{\prime}} - \varphi_ {\mu^{\prime}}) \| _{0} ^{2}
\\
&= C ^{2} C _{1} C _{2} r ^{2j-2k} \| \varphi_ {\nu^{\prime}} - \varphi_ {\mu^{\prime}} \| _{k} ^{2} \leqq 4 C ^{2} C _{1} C _{2} r ^{2j-2k}
\end{aligned}
\]
with a constant $C_2$. 

The last term converges, by $j < k$, to $0$ as $r \to \infty$. Therefore, $\lim\limits_{\nu, \mu \to \infty} \| T \varphi_{\nu'} - T \varphi_{\mu'} \|_k = 0$. 
\end{proof}


\yossection{4}{Schauder's Theorem}
\label{sec:X.4}

\begin{theoremn}[Schauder] An operator $T \in L(X, Y)$ is compact iff its dual operator $T'$ is compact.

\end{theoremn}
\begin{proof}
 Let $S, S'$ be the closed unit sphere in $X, Y'$, respectively. Let $T \in L(X, Y)$ be compact. Let $\{y_j'\}$ be an arbitrary sequence in $S'$. The functions $F_j(y) = \langle y, y_j' \rangle$ are equi-continuous in the sense that 
\[
\left| F _{j} (y) - F _{j} (z) \right| = \left| \langle y - z, y _{j} ^{\prime} \rangle \right| \leqq \left\| y - z \right\|. \]

Moreover, $\{F_j(y)\}$ is equi-bounded in $j$ on any bounded set of $y$, since $|F_j(y)| \leqq \|y\|$. Therefore, by the Ascoli-Arzelà theorem, as applied to the functions $\{F_j(y)\}$ defined on the compact set $(T\cdot S)^a$, we see that some subsequence $\{F_{j'}(y)\}$ converges uniformly in $y\in (T\cdot S)^a$. Hence $\langle Tx,y_j'\rangle = \langle x,T'y_j'\rangle$ converges uniformly in $x\in S$, and so $\{T'\cdot y_j'\}$ converges in the strong topology of $X'$. This proves that $T'$ is compact. 

Conversely, let $T'$ be compact. Then, by what we have proved above, $T''$ is compact. Hence, if $S''$ is the closed unit sphere in $X''$, $(T'' \cdot S'')$ is relatively compact. We know that $Y$ is isometrically embedded in $Y''$ (see Theorem~\ref{thm:IV.8.Theorem.2} in Chapter~\ref{chap:IV}, \ref{sec:IV.8}). Hence $T \cdot S \subseteq T'' \cdot S''$ and so $T \cdot S$ is relatively compact in the strong topology of $Y''$ and so in the strong topology of $Y$. Therefore, $T$ is compact. 
\end{proof}


\yossection{5}{The Riesz-Schauder Theory}
\label{sec:X.5}

\begin{lemman}[F. \textsc{Riesz} \cite{rickart273}] Let $V$ be a compact operator $\in L(X, X)$, where $X$ is a $B$-space. Then, for any complex number $\lambda_{0} \neq 0$, the range $R(\lambda_{0}I - V)$ is strongly closed.

\end{lemman}
\begin{proof}
 We may assume that $\lambda_{0}=1$. Let $\{x_{n}\}$ be a sequence of $X$ such that $y_{n}=(I-V)x_{n}$ converges strongly to $y$. If $\{x_{n}\}$ is bounded, then by, the compactness of the operator $V$, there exists a subsequence $\{x_{n'}\}$ such that $\{Vx_{n}\}$ converges strongly. Since $x_{n'}=y_{n'}+Vx_{n'},\{x_{n'}\}$ converges to some $x$, and so $y=(I-V)x$. 

We next assume that $\{\| x_n \| \}$ is unbounded. Set $T = (I - V)$ and put $\alpha_n = \text{ dis }(x_n, N(T))$, where $N(T) = \{x; Tx = 0\}$. Take a $w_n \in N(T)$ such that $\alpha_n \leqq \| x_n - w_n \| \leqq (1 + n^{-1}) \alpha_n$. Then $T(x_n - w_n) = Tx_n$, and so in the case when $\{\alpha_n\}$ is bounded, we can prove, as above, that $y \in R(T) = R(I - V)$. Suppose that $\lim\limits_{n \to \infty} \alpha_n = \infty$. Since $z_n = (x_n - w_n) / \| x_n - w_n \| $ satisfies $ \| z_n \| = 1$ and $\slim\limits_{n \to \infty} T z_n = 0$, we can prove, as above, that there exists a subsequence $\{z_{n'}\}$ such that $\slim\limits_{n \to \infty} z_{n'} = w_0$, $\slim\limits_{n \to \infty} T z_{n'} = 0$. Hence $w_0 \in N(T)$. But, if we put $z_n - w_0 = u_n$, then, in 
\[
x _{n} - w _{n} - w _{0} \left\| x _{n} - w _{n} \right\| = u _{n} \left\| x _{n} - w _{n} \right\|,
\]
the second and the third terms on the left belong to $N(T)$ so that we must have $\| u_n \| \cdot \| x_n - w_n \| \geqq \alpha_n$. This is a contradiction, since $s \text{- } \lim\limits_{n \to \infty} u_n = 0$, $\| x_n - w_n \| \leqq (1 + n^{-1}) \alpha_n$ and $\lim\limits_{n \to \infty} \alpha_n = \infty$. 

We are now able to prove the Riesz-Schauder theory; for convenience sake, we shall state the theory in a series of three theorems. 
\end{proof}
\begin{theorem}\label{thm:X.5.Theorem.1} Let $V$ be a compact operator $\in L(X, X)$. If $\lambda_0 \neq 0$ is not an eigenvalue of $V$, then $\lambda_0$ is in the resolvent set of $V$. 
\end{theorem}

\begin{proof}
 By the preceding Lemma and the hypothesis, the operator $T_{\lambda_{0}} = (\lambda_{0}I - V)$ gives a one-one map of $X$ onto the set $R(T_{\lambda_{0}})$ which is strongly closed in $X$. Hence, by the Corollary of the open mapping theorem in Chapter~\ref{chap:II}, \ref{sec:II.5}, $T_{\lambda_{0}}$ has a continuous inverse. We have to show that 

$R(T_{\lambda_0}) = X$. If not, the topological image $X_1 = T_{\lambda_0}X$ of $X$ is a proper closed subspace of $X$. Hence, if we set $X_2 = T_{\lambda_0}X_1$, $X_3 = T_{\lambda_0}X_2, \ldots$, then $X_{n+1}$ is a proper closed subspace of $X_n$ ( $n = 0, 1, 2, \ldots; X_0 = X$). By F. \textsc{Riesz}' theorem in Chapter~\ref{chap:III}, \ref{sec:III.2}, there exists a sequence $\{y_n\}$ such that $y_n \in X_n$, $\|y_n\| = 1$ and $\text{ dis }(y_n, X_{n+1}) \geqq 1/2$. Thus, if $n > m$, 
\[
\lambda_{0} ^{-1} \left(V y _{m} - V y _{n}\right) = y _{m} + \left\{- y _{n} - \left(T _ {\lambda_{0}} y _{m} - T _ {\lambda_{0}} y _{n}\right) / \lambda_{0} \right\} = y _{m} - y
\]
with some $y \in X_{m+1}$. 

Hence $\|Vy_n - Vy_m\| \geqq |\lambda_0| / 2$, contrary to the compactness of the operator $V$. 
\end{proof}

\begin{theorem}\label{thm:X.5.Theorem.2} Let $V$ be a compact operator $\in L(X, X)$. Then, \begin{enumerate}[(i)]
\item its spectrum consists of an at most countable set of points of the complex plane which has no point of accumulation except possibly $\lambda = 0$;
\item every non-zero number in the spectrum of $V$ is an eigenvalue of $V$ of finite multiplicity;
\item a non-zero number is an eigenvalue of $V$ iff it is an eigenvalue of $V'$.
\end{enumerate}
\end{theorem}

\begin{proof}
 By Theorem~\ref{thm:X.5.Theorem.1}, a non-zero number in the spectrum of $V$ is an eigenvalue of $V$. The same is also true of $V'$, since, by Schauder's theorem, $V'$ is compact when $V$ is. But the resolvent sets are the same for $V$ and $V'$ (see Chapter~\ref{chap:VIII}, \ref{sec:VIII.6}). Hence (iii) is proved. Since the eigenvectors belonging respectively to different eigenvalues of $V$ are linearly independent, the proofs of:
 (i) and (ii) are completed if we derive a contradiction from the following situation:
There exists a sequence $\{x _{n} \}$ of linearly independent vectors such that $V x _{n} = \lambda_{n} x (n = 1, 2, \dots)$ and $\lim\limits _{n \to \infty} \lambda_{n} = \lambda \neq 0.$


To derive a contradiction, we consider the closed subspace $X_{n}$ spanned by $x_{1}, x_{2}, \ldots, x_{n}$. By F. \textsc{Riesz}' theorem in Chapter~\ref{chap:III}, \ref{sec:III.2}, there exists a sequence $\{y_{n}\}$ such that $y_{n} \in X_{n}$, $\| y_{n} \| = 1$ and $\text{ dis }(y_{n}, X_{n-1}) \geqq 1/2$ ( $n = 2, 3, \ldots$). If $n > m$, then 
\[
\begin{aligned}
&\lambda_{n} ^{-1} V y _{n} - \lambda_{m} ^{-1} V y _{m} = y _{n} + (- y _{m} - \lambda_{n} ^{-1} T _ {\lambda_{n}} y _{n} + \lambda_{m} ^{-1} T _ {\lambda_{m}} y _{m}) = y _{n} - z, \text{ where } z \in X _{n-1}.
\end{aligned}
\]

For, if $y_{n} = \sum\limits_{j=1}^{n} \beta_{j} x_{j}$, then we have $y_{n} - \lambda_{n}^{-1} V y_{n} = \sum\limits_{j=1}^{n} \beta_{j} x_{j} - \sum\limits_{j=1}^{n} \beta_{j} \lambda_{n}^{-1} \lambda_{j} x_{j} \in X_{n-1}$ and similarly $T_{\lambda_m} y_m \in X_m$. Therefore $\|\lambda_n^{-1} V y_n - \lambda_m^{-1} V y_m\| \geqq 1/2$. This contradicts the compactness of $V$ combined with the hypothesis $\lim\limits_{n \to \infty} \lambda_n \neq 0$. 
\end{proof}
\begin{theorem}\label{thm:X.5.Theorem.3} Let $\lambda_0 \neq 0$ be an eigenvalue of a compact operator $V \in L(X, X)$. Then $\lambda_0$ is also an eigenvalue of $V'$ by the preceding theorem. We can prove: \begin{enumerate}[(i)]
\item the multiplicities for the eigenvalue $\lambda_0$ are the same for $V$ and $V'$.
\item The equation $(\lambda_0 I - V)x = y$ admits a solution $x$ iff $y \in N(\lambda_0 I' - V')^\perp$, that is, iff $V'f = \lambda_0 f$ implies $\langle y, f \rangle = 0$.
\item The equation $(\lambda_0 I' - V'f) = g$ admits a solution $f$ iff $g \in N(\lambda_0 I - V)^\perp$, that is, iff $Vx = \lambda_0 x$ implies $\langle x, g \rangle = 0$.
\end{enumerate}
\end{theorem}

\begin{proof}
 Since the eigenvalue $\lambda_{0} \neq 0$ is an isolated singularity of the resolvent $R(\lambda; V) = (\lambda I - V)^{-1}$, we can expand $R(\lambda; V)$ in Laurent series 
\[
R (\lambda ; V) = \sum\limits_{n = - \infty} ^{\infty} (\lambda - \lambda_{0}) ^{n} A _{n}. \]

We are particularly interested in the residue $A_{-1} = (2\pi i)^{-1} \int_{|\lambda - \lambda_0| = \varepsilon} R(\lambda; V) \dd\lambda$. As was proved in Chapter~\ref{chap:VIII}, \ref{sec:VIII.8}, $A_{-1}$ is an idempotent, i.e., $A_{-1}^2 = A_{-1}$. If we set $(\lambda I - V)^{-1} = \lambda^{-1}I + V_\lambda$, then from $(\lambda I - V)(\lambda^{-1}I + V_\lambda) = I$ we obtain $V_\lambda = V(\lambda^{-1}V_\lambda + \lambda^{-2}I)$, and so $V_\lambda$ is compact when $V$ is. Hence, by 
\[
\begin{aligned}
A _{-1} &= (2 \pi i) ^{-1} \int_ {| \lambda - \lambda_{0} | = \varepsilon} R (\lambda ; V) \dd \lambda = (2 \pi i) ^{-1} \int_ {| \lambda - \lambda_{0} | = \varepsilon} \lambda^{-1} \dd \lambda \cdot I
\\
+ (2 \pi i) ^{-1} \int_ {| \lambda - \lambda_{0} | = \varepsilon} V _{\lambda} \dd \lambda &= (2 \pi i) ^{-1} \int_ {| \lambda - \lambda_{0} | = \varepsilon} V _{\lambda} \dd \lambda .
\end{aligned}
\]

Thus, by the Theorem in Chapter~\ref{chap:X}, \ref{sec:X.2}, $A_{-1}$ is a compact operator. 

Therefore, by $A_{-1}X = A_{-1}(A_{-1}X)$ and the compactness of $A_{-1}$, the unit sphere of the normed linear space $A_{-1}X$ is relatively compact. Hence, by F. \textsc{Riesz}' theorem in Chapter~\ref{chap:III}, \ref{sec:III.2}, the range $R(A_{-1})$ is of finite dimension. On the other hand, $Vx = \lambda_0x$, $x \neq 0$, implies that $(\lambda I - V)^{-1}x = (\lambda - \lambda_0)^{-1}x$ by $(\lambda I - V)x = (\lambda - \lambda_0)x$, and so $A_{-1}x = (2\pi i)^{-1}\int_{|\lambda - \lambda_0| = \varepsilon}(\lambda - \lambda_0)^{-1}\dd\lambda \cdot x = x$. Therefore, the eigenvalue equation $Vx = \lambda_0x$ is equivalent to $Vx = \lambda_0x$, $x \in R(A_{-1})$. In the same way, we can prove that the eigenvalue equation $V'f = \lambda_0f$ is equivalent to $V'f = \lambda_0f$, $f \in R(A_{-1}')$. But $R(A_{-1})$ and $R(A_{-1}')$ are of the same dimension. For, $A_{-1}'f = g$ satisfies $A_{-1}'g = A_{-1}'(A_{-1}'f) = g$, and this is equivalent to $\langle x, g \rangle = \langle A_{-1}x, g \rangle$ for all $x \in X$ and so the functional $g$ may be considered as a functional defined on the finite dimensional space $R(A_{-1})$. 

Now, by the well-known theorem in matrix theory, the eigenvalue equation $Vx = \lambda_{0}x$ (in $R(A_{-1})$) and its transposed equation $V'f = \lambda_{0}f$ (in $R(A_{-1}')$) both have the same number of linearly independent solutions. We have thus proved:
 (i). The propositions (ii) and (iii) are already proved by the Lemma and the closed range theorem (Chapter~\ref{chap:VII}, \ref{sec:VII.5}).
\end{proof}

Extension of the Riesz-Schauder Theory. Let a power $V^{n}$ of $V \in L(X, X)$ be compact for some positive integer $n$. Then, by the spectral mapping theorem in Chapter~\ref{chap:VIII}, \ref{sec:VIII.7}, $\sigma(V^{n}) = \sigma(V)^{n}$ and, by the compactness of $V^{n}$, $\sigma(V^{n})$ is either a finite set or else a countable set accumulating only at $0$. Therefore, $\sigma(V)$ is either a finite set or a countable set accumulating 
only at $0$. $V^n$ being compact, 
\[
(2 \pi i) ^{-1} \int_ {| \lambda - \lambda_{0} | = \varepsilon} R (\lambda ; V ^{n}) \dd \lambda
\]
is, for any $\lambda_0 \neq 0$ of $\sigma(V^n)$ and sufficiently small $\varepsilon > 0$, of finite-dimensional range. Hence $\lambda_0$ is a pole of $R(\lambda; V^n)$ (see Chapter~\ref{chap:VIII}, \ref{sec:VIII.8}). But $(\lambda^n I - V^n) = (\lambda - V)(\lambda^{n-1}I + \lambda^{n-2}V + \cdots + V^{n-1})$ and hence 
\[
(\lambda^{n} I - V ^{n}) ^{-1} (\lambda^{n-1} I + \dots + V ^{n-1}) = (\lambda I - V) ^{-1},
\]
which proves that any $\lambda_0 \neq 0$ of $\sigma(V^n)$ is a pole of $R(\lambda; V)$ and so is an eigenvalue of $V$. These facts enable us to extend the Riesz-Schauder theory to operators $V$ for which some power $V^n$ is compact. This extension is highly important in view of its application to concrete problems of integral equations, such as the Dirichlet problem pertaining to potentials. See, e.g., O. D. \textsc{Kellogg} \cite{kellogg148}. It can be proved that the Riesz-Schauder theory for $\lambda_0 = 1$ is valid also for an operator $V \in L(X, X)$ if there exist a positive integer $m$ and a compact operator $K \in L(X, X)$ such that $\|K - V^m\| < 1$. See K. \textsc{Yosida} \cite{yosida352}. It is to be noted here that, if $K_1(s, t)$ and $K_2(s, t)$ are bounded measurable for $0 \leqq s$, $t \leqq 1$, then the integral operator $T$ defined by 
\[
x (s) \rightarrow (T x) (s) = \left(K _{1} K _{2} x\right) (s), \text{ where } \left(K _{j} x\right) (s) = \int_{0} ^{1} K _{j} (s, t) x (t) \dd t,
\]
is compact as an operator $L(L^{1}(0,1),L^{1}(0,1))$. See K. \textsc{Yosida}-Y. \textsc{Mimura}-S. \textsc{Kakutani} \cite{yosida353}. 


\yossection{6}{Dirichlet's Problem}
\label{sec:X.6}

Let $G$ be a bounded open domain of $\mathbb{R}^{n}$, and 
\[
L = \sum\limits_{| s |, | t | \leqq m} D ^{s} c _{st} (x) D ^{t}
\]
a strongly elliptic differential operator with real $C^\infty(G^a)$ coefficients $c_{st}(x) = c_{ts}(x)$. We shall deal only with real-valued functions. Let $f \in L^2(G)$ and $u_1 \in H^m(G)$ be given. Consider a distribution solution $u_0 \in L^2(G)$ of 
\begin{equation}\label{eq:X.6.1}
L u = f \text{ such that } (u _{0} - u _{1}) \in H _{0} ^{m} (G).
\end{equation}

The latter condition $(u_0 - u_1) \in H_0^m(G)$ means that each of the distributional derivatives 
\begin{equation}\label{eq:X.6.2}
\left(D ^{j} u _{0} - D ^{j} u _{1}\right) \text{ for } | j | \leqq m
\end{equation}
is the $L^{2}(G)$-limit of a sequence $\{D^{j}\varphi_{h,j}\}$, where $\varphi_{h,j} \in C_{0}^{\infty}(G)$ (see Chapter~\ref{chap:I}, \ref{sec:I.10}). Thus it gives roughly the boundary conditions: 
\begin{equation}\label{eq:X.6.3}
D ^{j} u _{0} = D ^{j} u _{1} \text{ on the boundary } \partial G \text{ of } G \text{ for } | j | < m.
\end{equation}

In such a sense, \eqref{eq:X.6.1} will be called a Dirichlet's problem for the operator 

$L$. We shall follow the treatment of the problem as formulated and solved by L. \textsc{Garding} \cite{garding74}. 

We first solve 
\begin{equation}\label{eq:X.6.4}
u + \alpha L u = f, (u - u _{1}) \in H _{0} ^{m} (G),
\end{equation}
where the positive constant $\alpha$ is so chosen that Gårding's inequality 
\begin{equation}\label{eq:X.6.5}
(\varphi + \alpha L ^{*} \varphi , \varphi) _{0} \geqq \delta \| \varphi \| _{m} ^{2} \text{ holds whenever } \varphi \in C _{0} ^{\infty} (G)
\end{equation}

Here $L^{*} = \sum\limits_{|s|, |t| \leqq m} (-1)^{|s|+|t|} D^{t} c_{st}(x) D^{s}$ and $\delta$ is a positive constant. The existence of such an $\alpha$ is guaranteed if the coefficients $c_{st}(x)$ are continuous on the closure $G^{a}$ of $G$. We also have, by $m$-times partial differentiation, the inequality 
\begin{equation}\label{eq:X.6.6}
\left| \left(\varphi + \alpha L ^{*} \varphi , \psi\right) _{0} \right| \leqq \gamma \left\| \varphi \right\| _{m} \cdot \left\| \psi \right\| _{m} \text{ whenever } \varphi , \psi \in C _{0} ^{\infty} (G),
\end{equation}
where $\gamma$ is another positive constant independent of $\varphi$ and $\psi$. 

We have, for $u_{1} \in H^{m}(G)$ and $\varphi \in C_0^\infty(G)$, 
\[
(L ^{*} \varphi , u _{1}) _{0} = \sum\limits_{s,t} ((- 1) ^{|s|+|t|} D ^{t} c _{st} D ^{s} \varphi , u _{1}) _{0} = \sum\limits_{s,t} (- 1) ^{|s|} (c _{st} D ^{s} \varphi , D ^{t} u _{1}) _{0},
\]
by partial differentiation. By Schwarz' inequality, we obtain, remembering that the coefficients $c_{st}$ are bounded on $G^{a}$, 
\[
\left| \left(L ^{*} \varphi , u _{1}\right) _{0} \right| \leqq \eta \sum\limits_{| s |, | t | \leqq m} \left\| D ^{s} \varphi \right\| _{0} \left\| D ^{t} u _{1} \right\| _{0} \quad \left(\sup\limits _{s,t;x} \left| c _{st} (x) \right| = \eta\right).
\]

The right hand side is smaller than constant times $\|\varphi\|_m$. 

Thus the linear functional 
\[
F (\varphi) = \left(\varphi + \alpha L ^{*} \varphi , u _{1}\right) _{0}, \varphi \in C _{0} ^{\infty} (G),
\]
can be extended to a bounded linear functional defined on $H_{0}^{m}(G)$ which is the completion of $C_{0}^{\infty}(G)$ with respect to the norm $\left\|\varphi\right\|_{m}$. Similarly, we see, from 
\[
\left| (\varphi , f) _{0} \right| \leqq \left\| \varphi \right\| _{0} \cdot \left\| f \right\| _{0} \leqq \left\| \varphi \right\| _{m} \cdot \left\| f \right\| _{0},
\]
that the linear functional $(\varphi, f)_0$ of $\varphi \in C_0^\infty(G)$ can be extended to a bounded linear functional of $\varphi \in H_0^m(G)$. Hence, by F. \textsc{Riesz}' representation theorem as applied to the Hilbert space $H_0^m(G)$, there exists an $f' = f'(f, u_1) \in H_0^m(G)$ such that 
\[
(\varphi , f) _{0} - (\varphi + \alpha L ^{*} \varphi , u _{1}) _{0} = (\varphi , f ^{\prime}) _{m} \text{ whenever } \varphi \in C _{0} ^{\infty} (G).
\]

Hence, by the Milgram-Lax theorem in Chapter~\ref{chap:III}, \ref{sec:III.7}, applied to the Hilbert space $H_0^m(G)$, we have 
\begin{equation}\label{eq:X.6.7}
(\varphi , f) _{0} - (\varphi + \alpha L ^{*} \varphi , u _{1}) _{0} = (\varphi , f ^{\prime}) _{m} = B (\varphi , S f ^{\prime}), S f ^{\prime} \in H _{0} ^{m} (G),
\end{equation}
where 
\begin{equation}\label{eq:X.6.8}
B (\varphi , \psi) = (\varphi + \alpha L ^{*} \varphi , \psi) _{0} \quad \text{ for } \quad \varphi \in C _{0} ^{\infty} (G), \psi \in H _{0} ^{m} (G).
\end{equation}

Thus 
\[
(\varphi , f) _{0} = (\varphi + \alpha L ^{*} \varphi , u _{1} + S f ^{\prime}) _{0} \quad \text{ whenever } \quad \varphi \in C _{0} ^{\infty} (G),
\]
and so $u_0 = u_1 + Sf'$ is the desired solution $\in L^2(G)$ of \eqref{eq:X.6.4}. 

We shall next discuss the original equation \eqref{eq:X.6.1}. If $u_0 \in L^2(G)$ satisfies \eqref{eq:X.6.1}, then $u_2 = u_0 - u_1 \in H_0^m(G)$ satisfies 
\[
(u _{0}, L ^{*} \varphi) _{0} = (u _{1}, L ^{*} \varphi) _{0} + (u _{2}, L ^{*} \varphi) _{0} = (f, \varphi) _{0}, \varphi \in C _{0} ^{\infty} (G).
\]

We obtain, by partial integration as above, 
$\left|(u_1, L^*\varphi)_0\right| \leqq a$ constant times $\left\|\varphi\right\|_m$, 
\[
\left| (f, \varphi) _{0} \right| \leqq \left\| f \| _{0} \right. \left\| \varphi \| _{0} \right. \leqq \left\| f \| _{0} \cdot \right. \left\| \varphi \| _{m} \cdot \right.
\]

We may thus apply F. \textsc{Riesz}' representation theorem in $H_0^m(G)$ to the linear functional $(f, \varphi)_0 - (u_1, L^*\varphi)_0$ of $\varphi \in C_0^\infty(G)$. Hence there exists a uniquely determined $v \in H_0^m(G)$ such that 
\[
(f, \varphi) _{0} - (u _{1}, L ^{*} \varphi) _{0} = (v, \varphi) _{m} \quad \text{ whenever } \quad \varphi \in C _{0} ^{\infty} (G).
\]

By the Milgram-Lax theorem, applied to $(v,\varphi)_m$, we obtain an $S_1v\in H_0^m (G)$ such that 
\[
(v, \varphi) _{m} = B (S _{1} v, \varphi) \quad \text{ whenever } \quad \varphi \in C _{0} ^{\infty} (G), v \in H _{0} ^{m} (G).
\]

Thus the Dirichlet problem \eqref{eq:X.6.1} is equivalent to the problem: For a given $S_{1}v\in H_{0}^{m}(G)$, find a solution $\pmb{u}_{2}\in H_{0}^{m}(G)$ of 
\begin{equation}\tag{1'}\label{eq:X.6.1'}
(u _{2}, L ^{*} \varphi) _{0} = B (S _{1} v, \varphi), \varphi \in C _{0} ^{\infty} (G).
\end{equation}

Now, for a given $u \in L^2(G) = H_0^0(G)$, 
\[
\left| (u, \varphi) _{0} \right| \leqq \left\| u \| _{0} \cdot \| \varphi \| _{0} \right. \leqq \left\| u \| _{0} \cdot \| \varphi \| _{m} \right.
\]
so that, by F. \textsc{Riesz}' representation theorem in the Hilbert space $H_0^m(G)$, there exists a uniquely determined $u' = Tu \in H_0^m(G)$ such that, whenever $\varphi \in C_0^\infty(G)$, 
\[
(u, \varphi) _{0} = (u ^{\prime}, \varphi) _{m} \quad \text{ and } \quad \| u ^{\prime} \| _{m} \leqq \| u \| _{0}.
\]

Hence, by the Milgram-Lax theorem, we obtain 
\begin{equation}\label{eq:X.6.9}
(u, \varphi) _{0} = (u ^{\prime}, \varphi) _{m} = B \left(S _{1} u ^{\prime}, \varphi\right) = B \left(S _{1} T u, \varphi\right), \| S _{1} T u \| _{m} \leqq \delta^{-1} \| u \| _{0}.
\end{equation}

Therefore, by \eqref{eq:X.6.1}, we have, whenever $\varphi \in C_0^\infty(G)$, 
\[
\begin{aligned}
&B (u _{2}, \varphi)  = (u _{2}, \varphi + \alpha L ^{*} \varphi) _{0} = (u _{2}, \varphi) _{0} + \alpha (u _{2}, L ^{*} \varphi) _{0} = B (S _{1} T u _{2}, \varphi) + \alpha B (S _{1} v, \varphi),
\end{aligned}
\]
that is, 
\[
B \left(u _{2} - S _{1} T u _{2} - \alpha S _{1} v, \varphi\right) = 0.
\]

Because of the positivity $B(\varphi, \varphi) > 0$ of $B$, we must have 
\begin{equation}\tag{1''}\label{eq:X.6.1''}
u _{2} - S _{1} T u _{2} = \alpha S _{1} v.
\end{equation}

The right hand term $\alpha S_1 v \in H_0^m(G)$ is a known function. By $\|S_1 T u\|_m \leqq \delta^{-1} \|u\|_0$, we see that the operator $S_1 T$ defined on $H_0^m(G)$ into $H_0^m(G)$ is compact (the Rellich-Gårding theorem in Chapter~\ref{chap:X}, \ref{sec:X.3}). Therefore, we may apply the Riesz-Schauder theory to the effect that one of the following alternatives holds: 

Either the homogeneous equation $u - S_1 Tu = 0$ has a non-trivial solution $u \in H_0^m(G)$, or the inhomogeneous equation $u - S_1 Tu = w$ has, for every given $w \in H_0^m(G)$, a uniquely determined solution $u \in H_0^m(G)$. 

The first alternative corresponds to the case $(u, \varphi + \alpha L^{*}\varphi)_{0} = (u, \varphi)_{0}$, that is, to the case $Lu = 0$. Hence, returning to the original equation \eqref{eq:X.6.1}, we have 

\begin{theoremu} One of the following alternatives holds: Either \begin{enumerate}[i)]
\item the homogeneous equation $L u = 0$ has a non-trivial solution $u \in H_0^m(G)$
\item for any $f \in L^2(G)$ and any $u_1 \in H^m(G)$, there exists a uniquely determined solution $u_0 \in L^2(G)$ of $L u = f$, $u - u_1 \in H_0^m(G)$.
\end{enumerate}
\end{theoremu}
